% Katalog: Einstieg, ID-System, Sektionsübersicht (Palette Slot 0).
\ZSFTitleHeader{Baustein-Katalog}{\ZSFOwnerDisplayName}
\StartFrontChapter[Kat]{Lesehilfe}

\begin{runintext}
  Jeder sichtbare Baustein des Templates einmal, mit identischem Mustertext.
  Was sich zwischen zwei Einträgen unterscheidet, ist deshalb der Baustein
  und nie der Inhalt.
\end{runintext}

\SubsectionBar*{Aufbau eines Eintrags}
\ZSFindex{Eintrag}

\begin{defbox}[B-00 · Aufbau]
  \Muster
\end{defbox}
\Zweck{Titel trägt ID und Bausteinnamen, der Inhalt ist überall derselbe,
diese Zeile nennt den Zweck.}

\begin{defbox}[Die drei Teile][weight=quiet]
  \ZSFlabel{ID:} adressierbar — »C-11 raus« genügt als Auftrag.
  \ZSFlabel{Name:} der Baustein oder der Regler samt Wert.
  \ZSFlabel{Zweckzeile:} wofür er gedacht war, in einem Satz.
\end{defbox}

\SubsectionBar*{Sektionen}
\ZSFindex{Sektionen}

\begin{tablebox}[Wo was steht]
  \begin{ZSFtable}{Z{0.34} Y{1.10} Y{1.56}}
    \ZSFheaderRow \ZSFhead{ID} & \ZSFhead{Kap.} & \ZSFhead{Inhalt} \\
    A & \ZSFsectionref{kat:struktur} & Kapitel, Balken, Umbruch \\
    B & \ZSFsectionref{kat:boxen} & die Box-Umgebungen \\
    C & \ZSFsectionref{kat:regler} & jeder Regler, jeder Wert \\
    D & \ZSFsectionref{kat:tabellen} & ZSFtable, valuegrid, Gruppe \\
    E & \ZSFsectionref{kat:marker} & Inline-Semantik, Verweise \\
    F & \ZSFsectionref{kat:bilder} & Abbildungen \\
    G & \ZSFsectionref{kat:formeln} & Formelsatz, Zielketten \\
    H & \ZSFsectionref{kat:module} & Opt-in-Module \\
    I & \ZSFsectionref{kat:inventar} & Liste: Art, ID, Codestelle
  \end{ZSFtable}
\end{tablebox}
\Zweck{A bis H zeigen, wie ein Baustein aussieht. I sagt, was er ist und in
welcher Datei er lebt — die beiden Fragen nach der Entscheidung.}

\begin{defbox}[Vollständigkeit][tone=neutral, weight=quiet]
  102 gezeigte Einträge in A--H, alle 104 Namen der stabilen API in I.
  Die 32 Namen ohne eigene Darstellung — Preambel- und Kapitel-Schalter,
  Spacing- und Umbruch-Makros — stehen dort mit einem Strich statt einer ID.
\end{defbox}
\Zweck{Nichts aus der stabilen API fehlt; was man nicht sehen kann, ist als
solches gekennzeichnet.}

\begin{warnbox}[Kein Prüfgegenstand]
  Dieses Dokument prüft nichts und wird von nichts geprüft. Die
  Vorführ-Pflicht der stabilen API hängt weiterhin an \ZSFlabel{main.tex};
  ein hier gestrichener Eintrag verschwindet nicht aus dem System.
\end{warnbox}
